\documentclass[phy1200]{handout}

\title{Atomic spectroscopy}

\begin{document}

\maketitle

Spectroscopy is what it is called when you are able to see the full spectrum of a certain beam of light. That is, when you can split a beam of light into all its colors and see how much of each color is in the beam. 

Around the end of the 19th century, people started noticing that gas discharge lamps (think neon signs) produce really strange spectra. You'll see the difference once you start looking around with a spectroscope! To use the handheld ones, point the slit at a light source. You will see the colors the beam is made of to the left. 

\section{Various spectra}
Point your spectroscope at some white light sources (e.g. the room lights or the sky). What do you see?

Point your spectroscope at a hydrogen gas discharge tube. How is this different? Now try a helium discharge tube. How are these two gas spectra different? For more on why atomic spectra are like this, see the textbook.\footnote{Open Stax University Physics Volume 3 chapter 6.4, especially the first section and figure 6.17. .}

Use the computer spectroscope to view the spectrum of the light coming off the discharge tube. Explain how to read the computer spectrum plots. 

Given the spectrum of helium and hydrogen, why do we perceive the light from hydrogen as a hot pink/purple and the light from helium as white?

Using the hydrogen spectrum, find the two strongest lines. Measure their wavelengths. These are the 3$\to$2 and 4$\to$2 atomic transitions (see figure 6.17). What is the value of the Rydberg constant? 


\section{report}
Things to explain in your report:



\end{document}